% ============================================================
%  colophon.tex — Matemagica
%  Da includere in main.tex prima di \end{document}
%  Uso: \input{colophon}
% ============================================================

\cleardoublepage
\thispagestyle{empty}

\vspace*{\fill}

{\small
    \begin{center}
        {\large\bfseries\color{blu} Matemagica}\\[4pt]
        {\normalsize Libro di supporto per l'apprendimento della matematica nelle scuole medie}\\[2pt]
        {\normalsize Ispirato al Piano di studio DECS, Canton Ticino}
    \end{center}

    \medskip
    \noindent\rule{\linewidth}{0.4pt}
    \medskip

    \noindent\textbf{Autore:} Francesco ed Emma Piraneo G.\\
    \textbf{Composizione e ottimizzazione:} Claude Sonnet\\
    \textbf{Lingua:} Italiano\\
    \textbf{Fascia scolastica:} Scuola media, anni 1–4 (età 11–14 anni)\\
    \textbf{Formato:} A4, ottimizzato per stampa e lettura digitale (PDF)\\
    \textbf{Composizione:} LuaLaTeX\\
    \textbf{Font testo:} Atkinson Hyperlegible — Braille Institute of America\\
    \textbf{Font matematica:} Latin Modern Math\\
    \textbf{Repository:} \href{https://github.com/fpiraneo/matemagica}{\texttt{github.com/fpiraneo/matemagica}}

    \medskip
    \noindent\rule{\linewidth}{0.4pt}
    \medskip

    \noindent\textbf{Licenza d'uso}

    \smallskip
    \noindent Quest'opera è distribuita con licenza\\
    \textbf{Creative Commons Attribuzione – Non Commerciale –}\\
    \textbf{Condividi allo stesso modo 4.0 Internazionale (CC BY-NC-SA 4.0).}

    \smallskip
    \noindent\href{https://creativecommons.org/licenses/by-nc-sa/4.0/deed.it}%
    {\texttt{creativecommons.org/licenses/by-nc-sa/4.0/deed.it}}

    \medskip

    \noindent\textbf{È vietato qualsiasi uso commerciale} di questo materiale, anche parziale.
    Ogni riproduzione o adattamento — anche parziale — deve includere
    un'attribuzione esplicita all'autore e alla fonte originale.

    \medskip

    \noindent Formula di attribuzione richiesta:

    \smallskip
    \begin{center}
        \textit{«Matemagica\\
            di Francesco ed Emma Piraneo Giuliano.\\
            Distribuito con licenza CC BY-NC-SA 4.0.\\
            Fonte: \href{https://github.com/fpiraneo/matemagica}{\texttt{github.com/fpiraneo/matemagica}}»}
    \end{center}

    \medskip
    \noindent\rule{\linewidth}{0.4pt}
    \medskip

    \noindent\textbf{Avvertenze}

    \smallskip
    \noindent I contenuti sono stati redatti con cura ispirandosi al Piano di studio DECS,
    ma non garantiscono completezza o aggiornamento rispetto a versioni successive
    del curricolo ufficiale. Il materiale è fornito \textit{«così com'è»}, senza garanzie
    di alcun tipo. L'autore declina ogni responsabilità per eventuali errori od omissioni.

    \smallskip
    \noindent Questo progetto non ha alcun endorsement ufficiale da parte del DECS,
    del Canton Ticino o di qualsiasi istituzione scolastica.
    Il materiale non sostituisce l'insegnamento frontale di un docente qualificato.

    \medskip
    \noindent\rule{\linewidth}{0.4pt}
    \medskip

    \begin{center}
        {\footnotesize Copyright \textcopyright\ Francesco ed Emma Piraneo Giuliano, \the\year.\\
            Tutti i diritti non espressamente concessi sono riservati.}
    \end{center}

} % fine \small

\vspace*{\fill}